% === Thread 4: Reasoning & Verification ===
\section{Thread 4: Reasoning \& Verification}
\label{sec:t4}

\begin{quote}
\textit{Core question: How can we enable models to perform reliable multi-step reasoning and verify intermediate steps?}
\end{quote}

\noindent
Reasoning is one of the most exciting capabilities of large language models and the most rapidly evolving direction in the post-training landscape.
A model that has undergone only pre-training, when confronted with mathematical or logical tasks requiring multi-step derivation,
tends to directly output an answer without showing any intermediate process, resulting in extremely high error rates and an inability to trace the source of errors.
The core challenge of the post-training stage is:
\textbf{How can we teach models to ``show their work,''
and enable them to self-verify each step while generating a reasoning chain?}

The evolutionary trajectory of this thread can be summarized in four stages:
\textbf{(1) Discovering the existence of reasoning capabilities} (CoT prompting demonstrates that few-shot examples can activate reasoning chains);
\textbf{(2) Improving reasoning reliability} (Self-Consistency and verification mechanisms);
\textbf{(3) Internalizing reasoning capabilities into model parameters} (bootstrapping methods such as STaR and ReST, as well as RL training with GRPO);
\textbf{(4) Deep integration of reasoning and computation} (test-time compute scaling, the o1/R1 paradigm).

The connections to other threads are especially close:
the idea of reward modeling flows from Thread~2 into this thread, giving rise to the process reward model
(\crossref{2}{sec:t2});
the long-chain reasoning capabilities developed in this thread provide core support for agent planning in Thread~6
(\crossref{6}{sec:t6});
and the efficiency techniques from Thread~7 have dramatically reduced the training and deployment costs of reasoning models
(\crossref{7}{sec:t7}).


% ============================================================
% 4.1 Problem Origin & Early Explorations
% ============================================================
\subsection{Problem Origin: The Discovery of Chain-of-Thought}
\label{sec:t4-origin}

\subsubsection{Scratchpads: The Earliest ``Scratch Paper''}

Before the term Chain-of-Thought was formally coined,
Nye et al.\ \cite{nye2021scratchpad} had already discovered a key phenomenon:
when models are allowed to write down intermediate computation steps on a ``scratchpad'' before producing the final answer,
their accuracy on tasks such as arithmetic and program execution tracing improves significantly.
Although simple, this work revealed a profound insight:
\textbf{the computational complexity achievable in a single forward pass of a Transformer is inherently limited},
and by generating intermediate tokens to ``extend'' the computation chain,
the model effectively increases its computation depth dynamically at inference time.

\subsubsection{Chain-of-Thought Prompting}
\label{sec:t4-cot}

\landmark{wei2022cot} systematized this discovery.
Wei et al.\ proposed \textbf{Chain-of-Thought (CoT) prompting}:
by including human-written reasoning processes in each few-shot prompt example
(rather than providing only input-output pairs),
model performance on multi-step reasoning tasks improves dramatically.

\begin{problembox}
Standard prompting provides only (question, answer) pairs:\\
\texttt{Q: Roger has 5 tennis balls...} $\rightarrow$ \texttt{A: 11}\\
The model is forced to complete all computation in a single forward pass, with no room for intermediate reasoning.
\end{problembox}

\begin{solutionbox}
CoT prompting inserts a reasoning chain before the answer:\\
\texttt{Q: Roger has 5 tennis balls...} $\rightarrow$
\texttt{A: Roger started with 5 balls. 2 cans of 3 = 6. 5 + 6 = 11. The answer is 11.}\\
The model progressively unfolds intermediate reasoning through autoregressive generation.
\end{solutionbox}

\paragraph{Key experimental results.}
The effectiveness of CoT prompting exhibits a pronounced \textbf{scale-dependent emergent} characteristic:
\begin{itemize}[leftmargin=2em]
  \item On small models ($<$10B parameters), CoT provides almost no benefit and may even degrade performance;
  \item On large models, the effect increases sharply with scale: PaLM-540B on GSM8K jumps from 17.9\% with standard prompting to 58.1\% with CoT prompting;
  \item On simpler arithmetic benchmarks such as MultiArith, PaLM-540B + CoT even reaches the level of prior task-specific fine-tuned models.
\end{itemize}

\noindent
The theoretical significance of this discovery is that
CoT does not ``teach'' the model to reason, but rather \textbf{activates} reasoning capabilities that were already acquired during pre-training but left unutilized.
This forms an interesting echo with the ``Superficial Alignment Hypothesis'' proposed by LIMA \cite{zhou2023lima} in Thread~1
(\crossref{1}{sec:t1})---
the core role of post-training may not be to inject new knowledge, but to activate the correct output patterns.

\subsubsection{Zero-Shot CoT}

While few-shot CoT requires humans to carefully design reasoning examples,
\landmark{kojima2022zeroshotcot} demonstrated a striking fact:
simply appending the phrase ``\textbf{Let's think step by step}'' to the end of a prompt
can activate the model's reasoning chain generation capability \textbf{without any examples}.

Kojima et al.\ termed this strategy \textbf{Zero-Shot CoT}.
On MultiArith, Zero-Shot CoT raised the accuracy of InstructGPT-175B from 17.7\% to 78.7\%;
its average performance across multiple reasoning benchmarks also surpasses standard zero-shot prompting.
This result indicates that large language models already possess the capability pattern for ``step-by-step thinking'' internally,
and only an extremely brief instruction is needed to ``unlock'' it.

\subsubsection{Least-to-Most Prompting}

Unlike CoT's ``all-at-once'' reasoning approach,
\citet{zhou2022least} proposed \textbf{Least-to-Most Prompting}---
\textbf{first decomposing complex problems into a sequence of simpler sub-problems},
then solving them progressively starting from the simplest, with each sub-problem's answer serving as context for subsequent sub-problems.
This compositional decomposition strategy performs particularly well on problems that require generalization beyond the complexity of the provided examples
(e.g., the SCAN compositional generalization benchmark),
with accuracy jumping from 16\% with CoT to 99.7\%.
Least-to-Most reveals an important design principle:
\textbf{the key to reasoning lies not only in ``thinking step by step,'' but also in ``decomposing the problem correctly''}.

\subsubsection{Tree of Thoughts}

The linear reasoning chain of CoT has inherent limitations when facing problems that require search and backtracking.
Yao et al.\ \cite{yao2023tot} proposed \textbf{Tree of Thoughts (ToT)},
which extends the reasoning process from a linear chain to a tree-structured search.
ToT generates multiple candidate ``thoughts'' at each step,
selects the most promising branches through an evaluation function (which can be the model's own self-evaluation),
and supports classic search strategies such as breadth-first search and depth-first search.
On tasks requiring creative search such as the Game of 24,
ToT achieves a success rate of 74\%, far surpassing CoT prompting at 4\%.

\paragraph{RAP: MCTS-driven reasoning.}
Hao et al.\ \cite{hao2023rap} proposed \textbf{Reasoning via Planning (RAP)},
which uses the LLM simultaneously as a \textbf{world model} (predicting the consequences of each reasoning step) and a \textbf{reasoning agent},
employing \textbf{Monte Carlo Tree Search (MCTS)} to replace the simpler BFS/DFS search in ToT.
The UCB exploration-exploitation balance of MCTS allows RAP to allocate search resources more efficiently under a limited computation budget.
This work combines classical planning algorithms with LLM reasoning,
and the application of MCTS was further developed in OmegaPRM \cite{luo2024omegaprm}'s automated process supervision
and rStar \cite{qi2024rstar}'s mutual reasoning.

\evolution{Chain-of-Thought: linear reasoning chains}{Tree of Thoughts / RAP: tree-structured search and planning}{need for backtracking and exploration}

\paragraph{PAL: Outsourcing computation to a code interpreter.}
Unlike the above methods that have the LLM perform all reasoning in natural language,
Gao et al.\ \cite{gao2022pal} proposed \textbf{Program-Aided Language models (PAL)}---
having the LLM generate the reasoning process as \textbf{executable programs} (rather than natural language reasoning chains),
which are then executed by a Python interpreter.
This division of labor---``reasoning for the language model, computation for the interpreter''---
raised PaLM-540B's accuracy on GSM8K from 56.5\% with CoT to \textbf{72.0\%},
with particularly notable advantages on problems requiring precise numerical computation.
PAL pioneered the \textbf{program-of-thought (PoT)} paradigm,
which was later absorbed as a core component of the hybrid CoT+PoT training strategy in works such as MAmmoTH \cite{yue2023mammoth},
and resonates with the tool-use ideas in Thread~6 (\crossref{6}{sec:t6}).

\subsubsection{CoT Distillation: Transferring Reasoning Capabilities to Smaller Models}

A key limitation of CoT is its scale-dependent nature---
small models cannot benefit from CoT prompting.
To address this, two parallel lines of research explored \textbf{CoT distillation}:

\begin{itemize}[leftmargin=2em]
  \item Magister et al.\ \cite{magister2023cotdistill} proposed using CoT reasoning chains generated by large models
    as SFT data to train smaller models,
    enabling a 0.3B model to approach the CoT prompting performance of a 540B model on arithmetic tasks.

  \item Hsieh et al.\ \cite{hsieh2023distillstep} proposed \textbf{Distilling Step-by-Step},
    which not only uses the large model's reasoning chains as labels,
    but also treats each step in the reasoning chain as an additional rationale supervision signal,
    achieving superior data efficiency compared to standard distillation.
    This method enabled a 770M T5 model to surpass the few-shot performance of the 540B PaLM on certain tasks using only 12.5\% of the training data.
\end{itemize}

\noindent
These distillation works laid the foundation for a core theme of Thread~4:
\textbf{how to internalize reasoning capabilities from inference-time prompt tricks into model parameters through post-training.}


% ============================================================
% 4.2 Foundation: Verification & Reward Modeling for Reasoning
% ============================================================
\subsection{Foundational Methods: Reasoning Verification and Reward Modeling}
\label{sec:t4-foundation}

Although CoT prompting activated reasoning chain generation capabilities,
the reliability of a single reasoning chain is far from sufficient to support high-stakes application scenarios.
A seemingly fluent reasoning chain may contain a subtle logical error at some intermediate step,
and this error propagates to all subsequent steps, leading to an incorrect final answer.
How to \textbf{verify} the correctness of reasoning chains became the second core problem of this thread.

\subsubsection{Self-Consistency: Multi-Path Voting}
\label{sec:t4-self-consistency}

\landmark{wang2023selfconsistency} proposed an intuitive and effective strategy:
rather than relying on a single greedy decoding reasoning chain,
generate multiple \textbf{diverse} reasoning paths through temperature sampling,
then perform \textbf{majority voting} on the final answers obtained from all paths.

\paragraph{Core method.}
Given a question $\mathbf{x}$, the decoding process of Self-Consistency is:
\begin{enumerate}[leftmargin=2em]
  \item Sample $m$ reasoning paths $(\mathbf{r}_1, a_1), \ldots, (\mathbf{r}_m, a_m)$ from the language model,
    where $\mathbf{r}_i$ is the reasoning chain and $a_i$ is the corresponding final answer;
  \item Select the answer with the highest frequency as the final prediction:
\end{enumerate}
\begin{equation}
\label{eq:self-consistency}
\hat{a} = \arg\max_{a} \sum_{i=1}^{m} \mathbf{1}(a_i = a)
\end{equation}

\noindent
This strategy can be understood as \textbf{marginalization} over the CoT reasoning process:
Wang et al.\ noted that Self-Consistency approximately implements
$\hat{a} = \arg\max_a \sum_{\mathbf{r}} P(\mathbf{r}, a \mid \mathbf{x})$,
i.e., selecting the answer with the highest posterior probability after summing over all reasoning paths.
Compared to greedy decoding or beam search, which only consider high-probability paths,
this ``sample-and-marginalize'' strategy better leverages the complementary information across diverse reasoning paths.

\paragraph{Key experimental results.}
Self-Consistency yields significant and consistent improvements across virtually all reasoning benchmarks:
\begin{itemize}[leftmargin=2em]
  \item GSM8K: GPT-3 code-davinci-002 improves from 60.1\% with CoT to 78.0\% (\textbf{+17.9\%});
  \item SVAMP: \textbf{+11.0\%}; AQuA: \textbf{+12.2\%};
  \item The method exhibits good robustness to sampling strategies (temperature, top-$k$);
  \item It is equally effective when combined with Zero-Shot CoT: GSM8K +26.2\%.
\end{itemize}

\noindent
The importance of Self-Consistency lies in revealing a key principle:
\textbf{investing more computational resources at inference time (generating more reasoning paths) can systematically improve reasoning accuracy}.
This principle became the direct precursor to subsequent research on test-time compute scaling (see \S\ref{sec:t4-convergence}).

\subsubsection{Self-Refine \& Reflexion: Iterative Self-Improvement}

Beyond the ``generate multiple paths + vote'' paradigm of Self-Consistency,
another line of research explored enabling models to \textbf{iteratively self-improve} their reasoning processes.

\begin{itemize}[leftmargin=2em]
  \item \textbf{Self-Refine} \cite{madaan2023selfrefine}:
    After generating an initial output, the model self-evaluates (critiques) the quality of its output,
    then modifies (refines) it based on the evaluation, iterating until a stopping condition is met.
    This approach demonstrated continuous improvement on tasks such as code optimization and dialogue response generation.

  \item \textbf{Reflexion} \cite{shinn2023reflexion}:
    Models the self-reflection process as a form of ``verbal reinforcement learning.''
    After completing a task, the agent generates a natural language reflection
    summarizing the causes of failure and strategies for improvement, storing it as episodic memory,
    which is then leveraged in subsequent trials to avoid repeating the same mistakes.
    Achieves 91\% pass@1 on HumanEval.
\end{itemize}

\paragraph{Limitations of self-correction.}
However, a systematic study by Huang et al.\ \cite{huang2024selfcorrect} showed that
\textbf{LLMs are actually unable to reliably self-correct their reasoning without external feedback}.
Across multiple reasoning benchmarks, self-correction even led to performance degradation.
Only when models receive external feedback signals (such as ground truth, tool outputs, or evaluations from stronger models)
can iterative improvement truly work.
This finding underscores the \textbf{irreplaceability of external verification mechanisms}.

\subsubsection{Outcome Reward Model (ORM) \& GSM8K}
\label{sec:t4-orm}

If Self-Consistency implicitly leverages the diversity of reasoning paths through majority voting,
then \textbf{training a dedicated verifier} to explicitly evaluate the quality of reasoning chains
represents a more refined verification strategy.

\landmark{cobbe2021gsm8k} made pioneering contributions in this direction.
Cobbe et al.\ presented two interrelated contributions:

\begin{enumerate}[leftmargin=2em]
  \item \textbf{The GSM8K dataset}: 8.5K high-quality grade school math word problems,
    each requiring 2--8 reasoning steps, with integer answers.
    This dataset became one of the most widely used benchmarks in subsequent reasoning research.

  \item \textbf{Outcome-based Verifier} (later termed Outcome Reward Model, ORM):
    Training a verifier model to judge whether the final answer of a reasoning chain is correct.
    At inference time, $N$ candidate reasoning chains are generated,
    the verifier scores each chain, and the highest-scoring answer is selected.
    Compared to majority voting, verifier-guided selection achieves higher accuracy for the same number of samples.
\end{enumerate}

\noindent
The idea behind ORM directly inherits the reward modeling framework from Thread~2
(\crossref{2}{sec:t2})---
training human judgments of reasoning quality into an automated scoring model.
However, a fundamental limitation of ORM is that
\textbf{it only evaluates the correctness of the final answer and cannot localize errors in intermediate reasoning steps}.
A chain that arrives at the correct answer through flawed reasoning (``right answer for wrong reasons'')
would be rated as correct by the ORM, which is unacceptable in scenarios requiring reliable reasoning.

\subsubsection{Process Reward Model (PRM): Let's Verify Step by Step}
\label{sec:t4-prm}

\landmark{lightman2023prm} directly addressed the above limitations of ORM.
Lightman et al.\ proposed the \textbf{Process Reward Model (PRM)}:
rather than evaluating only the final answer, it provides an independent correctness judgment for \textbf{each step} in the reasoning chain.

\paragraph{Core method.}
The PRM scores a reasoning chain $(s_1, s_2, \ldots, s_K)$ as follows:
\begin{equation}
\label{eq:prm-score}
\text{PRM-score}(\mathbf{s}) = \prod_{k=1}^{K} P(\text{correct} \mid s_1, \ldots, s_k)
\end{equation}
where each step $s_k$ is labeled as one of three categories: positive, negative, or neutral.
The training data for the PRM comes from step-by-step human annotations.

\paragraph{The PRM800K dataset.}
To train a high-quality PRM, Lightman et al.\ constructed the \textbf{PRM800K} dataset:
\begin{itemize}[leftmargin=2em]
  \item Approximately 800,000 step-level annotations covering 75,000 solution traces corresponding to 12,000 problems from the MATH dataset;
  \item Each step is labeled as positive (correct), negative (incorrect), or neutral (ambiguous);
  \item An \textbf{active learning} strategy was used to select the most informative samples for annotation,
    achieving a \textbf{2.6$\times$} improvement in annotation efficiency
    (i.e., the number of samples required to achieve equivalent performance under a fixed annotation budget was reduced to 1/2.6).
\end{itemize}

\paragraph{Key experimental results: PRM substantially outperforms ORM.}
In the best-of-$N$ evaluation on the MATH dataset ($N = 1860$):
\begin{itemize}[leftmargin=2em]
  \item PRM (best-of-$N$): \textbf{78.2\%}
  \item ORM (best-of-$N$): 72.4\%
  \item Majority Voting: 69.6\%
\end{itemize}

\noindent
PRM exceeds ORM by \textbf{5.8 percentage points} and majority voting by \textbf{8.6 percentage points}.
This result strongly demonstrates that \textbf{process-level supervision is superior to outcome-level supervision}.

\begin{openproblembox}
The annotation cost of PRM800K is extremely high (requiring domain experts in mathematics to check each step).
How to automate the annotation process for process supervision became a core open problem for subsequent research.
\end{openproblembox}

\subsubsection{Automating Process Supervision}

The tension between PRM's strong performance and its prohibitive annotation costs gave rise to a series of works on automating process supervision.

\paragraph{Math-Shepherd.}
Wang et al.\ \cite{wang2024mathshepherd} proposed \textbf{Math-Shepherd},
which uses an automatic annotation method based on a \textbf{completer model} to replace human annotation.
For each step $s_k$ in a reasoning chain,
Math-Shepherd has a model ``complete'' the remaining reasoning steps multiple times starting from $s_k$;
if the majority of completion paths reach the correct answer, the step is labeled as correct.
Without any human step-level annotations,
the PRM trained using this method matches or even surpasses the performance of human-annotated PRMs across multiple mathematical reasoning benchmarks.

\paragraph{OmegaPRM.}
Luo et al.\ \cite{luo2024omegaprm} proposed \textbf{OmegaPRM},
which employs Monte Carlo tree search (MCTS) to efficiently auto-annotate the correctness of reasoning steps.
Unlike Math-Shepherd's ``flat completion'' approach,
OmegaPRM uses tree search strategies to allocate computational resources more efficiently,
investing more sampling effort on critical steps that require fine-grained judgment
and quickly passing through simple steps, significantly improving annotation efficiency.

\paragraph{GenRM: Generative verification.}
Zhang et al.\ \cite{zhang2024genrm} proposed a different verification paradigm, \textbf{GenRM}
(Generative Reward Model):
rather than modeling verification as a discriminative classification problem (``this step is correct vs.\ incorrect''),
the verification model outputs its verification process and judgment through \textbf{natural language generation}.
GenRM leverages the language model's strengths in CoT reasoning,
unfolding a reasoning chain during the verification process to judge the correctness of each step,
and achieves superior out-of-distribution generalization compared to traditional discriminative PRMs.

\evolution{Human-annotated PRM (PRM800K)}{Automated PRM (Math-Shepherd, OmegaPRM)}{annotation cost bottleneck}


% ============================================================
% 4.3 RL for Reasoning: Bootstrapping & Policy Optimization
% ============================================================
\subsection{RL for Reasoning: Self-Improvement and Policy Optimization}
\label{sec:t4-rl}

The methods discussed in the previous two sections---CoT prompting, Self-Consistency, PRM---
primarily improve reasoning capabilities \textbf{at inference time} through prompt design or verifier-guided search.
This section turns to a more ambitious goal:
\textbf{internalizing reasoning capabilities into the model through post-training (modifying model parameters)},
so that the model can exhibit strong reasoning capabilities without relying on elaborate prompts or extensive sampling.

\subsubsection{STaR: Self-Taught Reasoner}
\label{sec:t4-star}

\landmark{zelikman2022star} proposed \textbf{STaR} (Self-Taught Reasoner),
a method for improving reasoning capabilities through \textbf{iterative self-training}.
The core idea of STaR is simple yet profound:
let the model generate its own reasoning chains, retain only those that arrive at the correct final answer as training data,
then fine-tune on these ``success cases,'' and iterate.

\paragraph{Algorithm procedure.}
Each iteration of STaR consists of the following steps:
\begin{enumerate}[leftmargin=2em]
  \item \textbf{Generate}: For each training problem $(x_i, y_i)$,
    have the current model $M_t$ generate a reasoning chain $\hat{r}_i$ and answer $\hat{y}_i$;
  \item \textbf{Filter}: Retain only samples where $\hat{y}_i = y_i$ (correct answers);
  \item \textbf{Rationalization}: For problems the model fails to answer correctly,
    provide the correct answer $y_i$ as a hint to the model,
    and have it ``reverse-engineer'' a reasoning chain that leads to the correct answer.
    These rationalized samples are also added to the training set;
  \item \textbf{Train}: Fine-tune the model on the collected (reasoning chain, correct answer) pairs
    to obtain $M_{t+1}$;
  \item Return to step 1 and repeat.
\end{enumerate}

\paragraph{The elegance of rationalization.}
The most critical innovation in STaR is the \textbf{rationalization} mechanism.
If only problems the model can already answer correctly are retained,
the training distribution becomes biased toward ``easy problems,'' preventing the model from making progress on harder ones.
Rationalization, by providing the correct answer as a hint,
enables the model to generate plausible reasoning chains even for problems it originally could not solve,
thereby pushing the ``learning frontier'' toward higher difficulty levels.

\paragraph{Deep connections to RL.}
Zelikman et al.\ proved that STaR can be understood as a special form of the REINFORCE algorithm.
Define the training objective as:
\begin{equation}
\label{eq:star-rl}
J(M, X, Y) = \sum_{(x,y) \in (X,Y)} \mathbb{E}_{\hat{r} \sim M(\cdot|x)}
\left[ \mathbf{1}(\hat{y} = y) \right]
\end{equation}
where $\hat{y}$ is the answer extracted from the reasoning chain $\hat{r}$.
Through the log-derivative trick, the gradient of this objective becomes:
\begin{equation}
\label{eq:star-gradient}
\nabla J = \sum_{(x,y)} \mathbb{E}_{\hat{r} \sim M}
\left[ \mathbf{1}(\hat{y} = y) \cdot \nabla \log P_M(\hat{r} \mid x) \right]
\end{equation}
This is identical in form to the REINFORCE gradient
(with reward $\mathbf{1}(\hat{y} = y)$ and policy $P_M(\hat{r} \mid x)$).
STaR's ``filter correct samples + SFT'' operation is effectively an approximate implementation of the REINFORCE gradient.

\paragraph{Key results.}
On CommonsenseQA,
STaR improves GPT-J 6B from an initial 36.6\% to 72.5\% through iterative self-training,
approaching the 73.0\% of GPT-3 175B few-shot fine-tuned (using only \textbf{1/30} of the parameters).
On GSM8K it reaches 10.7\%, which, although modest in absolute terms, demonstrates the viability of bootstrapping.

\subsubsection{V-STaR \& Quiet-STaR: Extensions}

\paragraph{V-STaR.}
Hosseini et al.\ \cite{hosseini2024vstar} proposed \textbf{V-STaR},
which introduces a \textbf{verifier} on top of the STaR framework.
The key innovation is that the incorrect reasoning chains discarded by STaR are not waste,
but valuable negative samples for training the verifier.
V-STaR simultaneously trains a generator (using correct samples for SFT) and a verifier (using correct/incorrect sample pairs) in each iteration,
with the two reinforcing each other: the generator produces more diverse reasoning chains,
while the verifier learns to more accurately distinguish correct from incorrect reasoning paths.

\paragraph{Quiet-STaR.}
Zelikman et al.\ \cite{zelikman2024quietstar} generalized the idea of STaR to
\textbf{the generation of all tokens}, not just explicit reasoning tasks.
Quiet-STaR has the model perform ``internal thinking'' before generating each token
(generating thought tokens that are not output),
then uses the hidden states of the thought tokens to improve the prediction of the next token.
The ambition of this work is to make reasoning the \textbf{default behavioral mode} of the language model,
rather than a special mode activated only when prompted.

\subsubsection{\texorpdfstring{ReST \& ReST\textsuperscript{EM}: Reinforced Self-Training}{ReST \& ReST-EM: Reinforced Self-Training}}

\paragraph{ReST.}
Gulcehre et al.\ \cite{gulcehre2023rest} proposed \textbf{ReST} (Reinforced Self-Training),
which is essentially a more systematic version of STaR.
ReST consists of two alternating steps:
\textbf{Grow} (generate samples using the current policy and filter them with a reward model) and
\textbf{Improve} (perform SFT on the filtered high-quality samples).
This framework can use more general reward signals than STaR, not limited to binary correctness.

\paragraph{ReST\textsuperscript{EM}.}
Singh et al.\ \cite{singh2024restem} reinterpreted ReST as an
\textbf{Expectation-Maximization (EM)} algorithm:
the E-step corresponds to ``generate and filter,'' and the M-step corresponds to ``train on filtered samples.''
ReST\textsuperscript{EM} systematically studied the scaling properties of this framework on reasoning tasks,
finding that on benchmarks such as MATH and HumanEval,
multiple iterations yield continued but diminishing performance gains (benefits drop significantly after 2--3 rounds),
and that using more generated samples (a larger candidate pool) and higher-quality filtering criteria lead to better final results.

\subsubsection{GRPO: Group Relative Policy Optimization}
\label{sec:t4-grpo}

\landmark{shao2024grpo} introduced
\textbf{Group Relative Policy Optimization (GRPO)} in DeepSeekMath,
representing a key simplification of the PPO framework from Thread~2.
The algorithmic details and unified perspective of GRPO have been derived in detail in \S\ref{sec:t2-grpo};
here we focus on its particular significance for reasoning tasks.

\paragraph{Why do reasoning tasks especially need GRPO?}
PPO faces two challenges on reasoning tasks that are more severe than in general dialogue scenarios:
\begin{enumerate}[leftmargin=2em]
  \item \textbf{The estimation dilemma of the value model}:
    The reward for mathematical reasoning is sparse and binary (the answer is either correct or incorrect),
    value signals for intermediate steps are extremely sparse,
    making value function estimation far more difficult than in general RLHF scenarios;
  \item \textbf{The necessity of large-scale sampling}:
    The accuracy on reasoning tasks is typically much lower than on general dialogue tasks,
    requiring a larger group size $G$ to ensure the existence of positive examples within the group,
    which strains the memory budget of PPO's four-model architecture.
\end{enumerate}

\noindent
GRPO addresses both problems by completely removing the value model
and replacing advantage estimation with group-relative reward normalization
(formulas in Eq.~\ref{eq:pref:grpo-advantage}--\ref{eq:pref:grpo-obj}).
On DeepSeekMath 7B, GRPO achieves 51.7\% on MATH (PPO: 50.1\%, SFT: 46.8\%)
and 88.2\% on GSM8K (PPO: 84.0\%),
delivering not only the best performance but also significantly reduced computational overhead by eliminating the value model.
This method later became the core algorithmic component in DeepSeek-R1's training.


% ============================================================
% 4.4 Convergence: Reasoning-Native Models
% ============================================================
\subsection{Convergence and Frontier: Reasoning-Native Models}
\label{sec:t4-convergence}

In the second half of 2024, all the preceding research threads---CoT, verification, RL for reasoning, test-time compute---
converged into an entirely new model paradigm: \textbf{reasoning-native models}.
These models acquire the ability to autonomously unfold long-chain reasoning at inference time
through large-scale RL training during the post-training stage,
no longer relying on manually designed CoT prompts or few-shot examples.

\subsubsection{OpenAI o1: The Turning Point}
\label{sec:t4-o1}

The release of \landmark{openai2024o1} marked the formal establishment of this paradigm.
Although OpenAI did not disclose the complete technical details,
its system card and subsequent technical reports revealed the following core elements:

\begin{itemize}[leftmargin=2em]
  \item \textbf{Large-scale RL training}:
    Using reinforcement learning in the post-training stage to teach the model
    to autonomously generate a long and detailed ``hidden chain of thought'' at inference time;
  \item \textbf{Hidden CoT}:
    Before answering, the model first generates an internal reasoning process that is not shown to the user,
    which can span thousands or even tens of thousands of tokens;
  \item \textbf{Test-time compute scaling}:
    o1's reasoning capability improves continuously with the computational resources
    (thinking time) invested at inference, exhibiting a novel type of scaling law.
\end{itemize}

\paragraph{Key results.}
\begin{itemize}[leftmargin=2em]
  \item AIME 2024 (American Invitational Mathematics Examination): 83.3\% (o1) vs 13.4\% (GPT-4o)
  \item Codeforces competition rating: 89\textsuperscript{th} percentile
  \item GPQA Diamond (graduate-level science questions): 78.0\%, surpassing domain experts
\end{itemize}

\noindent
The significance of o1 extends beyond specific benchmark scores.
It demonstrated a core proposition:
\textbf{post-training can teach a model not only ``how to answer'' (instruction following),
but also ``how to think'' (reasoning)}.
This opened an entirely new dimension for post-training research.

\subsubsection{DeepSeek-R1: The Emergence of Reasoning Capabilities}
\label{sec:t4-r1}

\landmark{deepseek2025r1} represents an important open-source replication and deepening of the o1 paradigm.
The greatest contribution of DeepSeek-R1 lies in revealing a striking phenomenon:
\textbf{reasoning capabilities can spontaneously emerge from pure RL training}.

\paragraph{R1-Zero: Pure RL reasoning without SFT.}
DeepSeek first conducted a bold experiment:
on DeepSeek-V3-Base (a model that had undergone only pre-training),
\textbf{without any SFT}, they directly applied GRPO for RL training.
The reward signal was extremely simple---consisting of only two rule-based components:
\begin{itemize}[leftmargin=2em]
  \item \textbf{Accuracy reward}: Whether the final answer matches the ground truth (binary, 0 or 1);
  \item \textbf{Format reward}: Whether the output follows the specified format
    (wrapping the reasoning process in \texttt{<think>...</think>} tags,
    placing the final answer in \texttt{<answer>...</answer>} tags).
\end{itemize}

\noindent
The results were striking. As training progressed, R1-Zero exhibited the following emergent behaviors:

\begin{enumerate}[leftmargin=2em]
  \item \textbf{Automatic lengthening of reasoning chains}:
    The model's average response length gradually grew from a few hundred tokens initially to several thousand tokens,
    indicating that the model autonomously learned to ``spend more time thinking'';

  \item \textbf{Emergence of self-reflection}:
    The model spontaneously produced self-reflective and backtracking behaviors in its reasoning chains,
    such as ``Wait, let me reconsider...'' and ``Hmm, that doesn't seem right...,''
    \textbf{none of which were explicitly encoded in the training data or reward};

  \item \textbf{``Aha moment''}:
    The DeepSeek team observed that the model suddenly ``had an epiphany'' at certain training stages---
    expressions like ``I found an interesting trick'' appeared in the reasoning chains,
    followed by a sudden leap in reasoning accuracy.
\end{enumerate}

\noindent
These emergent phenomena carry profound implications for AI research.
They suggest that \textbf{reasoning is not a skill that needs to be explicitly taught,
but rather a capability that can ``crystallize'' from pre-training through the right RL signals}.

\paragraph{Limitations of R1-Zero.}
Despite its impressive reasoning capabilities
(AIME 2024: 71.0\% pass@1, 86.7\% majority voting),
R1-Zero suffers from three main issues:
(1) poor readability---reasoning chains are a mixture of multiple languages with irregular formatting;
(2) divergence---reasoning chains sometimes extend indefinitely without converging to an answer;
(3) insufficient generality---suboptimal performance on non-mathematical/logical tasks.

\paragraph{The four-stage training pipeline of R1.}
To overcome the limitations of R1-Zero, DeepSeek designed a complete four-stage training pipeline for R1:

\begin{enumerate}[leftmargin=2em]
  \item \textbf{Stage 1: Cold-start SFT}\\
    Collecting several thousand high-quality long CoT examples (via few-shot prompting + human annotation)
    and performing SFT on the base model.
    The purpose is to provide the model with a good ``thinking format'' starting point,
    avoiding the formatting chaos observed in R1-Zero.

  \item \textbf{Stage 2: Reasoning-oriented RL}\\
    Conducting large-scale GRPO training on the cold-start SFT model,
    using rule-based rewards (accuracy + format),
    focusing on automatically verifiable tasks such as mathematics, coding, and logical reasoning.
    This stage involves massive training volume and is the core stage for the major improvement in reasoning capabilities.

  \item \textbf{Stage 3: Rejection sampling + SFT}\\
    Using the Stage 2 RL model to generate large quantities of reasoning chains,
    filtering high-quality samples through rejection sampling.
    Simultaneously mixing in high-quality SFT data for general tasks (writing, translation, question answering, etc.),
    and performing SFT on the combined dataset,
    enabling the model to maintain general capabilities while preserving reasoning ability.

  \item \textbf{Stage 4: All-scenario RL}\\
    The final stage of RL training uses \textbf{mixed rewards}:
    continuing to use rule-based rewards (accuracy) for reasoning tasks,
    and preference-based reward models for general tasks
    (\crossref{2}{sec:t2}),
    to simultaneously optimize helpfulness and harmlessness.
\end{enumerate}

\paragraph{Key results of R1.}
After the complete four-stage training, DeepSeek-R1 achieved the following performance:
\begin{itemize}[leftmargin=2em]
  \item AIME 2024: \textbf{79.8\%} pass@1, on par with OpenAI-o1-0912
  \item MATH-500: \textbf{97.3\%}
  \item Codeforces: 2029 Elo rating (surpassing 96.3\% of human contestants)
  \item GPQA Diamond: \textbf{71.5\%}
  \item Maintained strong general capabilities on open-ended dialogue and writing tasks
\end{itemize}

\paragraph{Distillation experiments.}
R1 also confirmed the distillability of reasoning capabilities:
using long CoT data generated by R1 to perform SFT on 1.5B--70B models from the Qwen-2.5 and Llama-3 families,
the distilled smaller models significantly outperformed standard models of equivalent size on reasoning benchmarks.
For example, DeepSeek-R1-Distill-Qwen-32B achieved 72.6\% on AIME 2024,
even surpassing OpenAI o1-mini (63.6\%).
This demonstrates that \textbf{reasoning capabilities acquired through large-model RL training
can be transferred to smaller models via straightforward SFT distillation},
echoing the CoT distillation work discussed in \S\ref{sec:t4-origin}.

\subsubsection{Kimi k1.5}

Moonshot AI's Kimi k1.5 \cite{moonshot2025kimi} validated the effectiveness of
the reasoning-native model paradigm from a different angle.
The unique contribution of k1.5 lies in its proposed \textbf{long2short} strategy:
first training a ``long thinking'' model (allowing extremely long reasoning chains),
then compressing reasoning efficiency through distillation and RL,
drastically shortening reasoning chain length while maintaining reasoning quality.
This strategy directly relates to the efficiency optimization of test-time compute.

\subsubsection{2025--2026: The Proliferation of Thinking Models}
\label{sec:t4-thinking-models}

2025 marked the transition of reasoning-native models from ``frontier experiments'' to ``industry default paradigm.''
After o1 and R1 demonstrated that reasoning capabilities can be acquired through RL,
the following year and a half of development proceeded along three main lines:
\textbf{(1) continued advances at the closed-source frontier} (o3 series),
\textbf{(2) explosive growth of open-source reasoning models} (DeepSeek's continued evolution, Qwen3),
\textbf{(3) systematic breakthroughs in RLVR training methodology} (DAPO and its successors).

\paragraph{The OpenAI o3 series.}
OpenAI released the successor series to o1 in 2025 \cite{openai2025o3}:
\begin{itemize}[leftmargin=2em]
  \item \textbf{o3} (April 2025): Achieved \textbf{88.9\%} on AIME 2025,
    a significant improvement over o1's 83.3\%;
    performance on the ARC-AGI benchmark was more than $3\times$ that of o1;
  \item \textbf{o4-mini} (April 2025): A lightweight reasoning model
    that dramatically reduces inference costs while maintaining strong reasoning capabilities;
  \item \textbf{o3-pro} (June 2025): A high-performance variant targeting complex scientific and engineering problems.
\end{itemize}
A key shift in the o3 series was the introduction of \textbf{native multimodal reasoning} and \textbf{agentic capabilities}---
the model is no longer just a text reasoner, but can also invoke tools, browse the web, and analyze images during the reasoning process
(\crossref{6}{sec:t6}; multimodal reasoning training in \crossref{5}{sec:t5-reasoning-rl}).
This marks the first deep convergence of the reasoning and agentic threads.

\paragraph{DeepSeek's post-R1 evolution.}
DeepSeek continued to iterate after R1:
\begin{itemize}[leftmargin=2em]
  \item \textbf{R1-0528} (May 2025): An enhanced version of R1
    that further scaled up RL training on the original architecture,
    achieving significant improvements across multiple reasoning benchmarks;
  \item \textbf{V3.1} (August 2025) \cite{deepseek2024v3}:
    Introduced a \textbf{thinking/non-thinking hybrid mode},
    where a single model supports both fast direct answering and deep reasoning modes,
    with users able to control whether reasoning chain generation is enabled via API parameters.
    This design philosophy later became an industry standard;
  \item \textbf{V3.2} (December 2025) \cite{deepseek2025v32}:
    Reasoning capability reached \textbf{IMO gold medal level},
    demonstrating performance on mathematical Olympiad-level problems approaching that of top human contestants.
\end{itemize}

\noindent
The evolution from R1 to V3.2 reveals an important trend:
\textbf{reasoning no longer exists as a separate model category},
but rather as a switchable mode within general-purpose models.
``Thinking models'' are becoming the default paradigm.

\paragraph{Qwen3 and QwQ.}
Alibaba's Qwen3 \cite{yang2025qwen3} contributed a systematic methodology for training reasoning models.
Its post-training adopts a \textbf{four-stage pipeline}:
\begin{enumerate}[leftmargin=2em]
  \item \textbf{Long CoT cold start}: SFT using curated long reasoning chain data;
  \item \textbf{Reasoning RL}: Large-scale GRPO/DAPO training on verifiable tasks;
  \item \textbf{Thinking mode fusion}: Fusing reasoning and non-reasoning modes within a single model,
    enabling the model to adaptively choose whether to activate reasoning chains based on problem complexity;
  \item \textbf{General RL}: All-scenario RL training with mixed rewards.
\end{enumerate}
The third stage, thinking mode fusion, is Qwen3's unique contribution---
it goes beyond DeepSeek V3.1's ``dual mode'' concept
by making reasoning mode switching an explicit training stage in the post-training pipeline.
QwQ-32B \cite{qwen2025qwq}, as the first reasoning model in the Qwen3 series,
achieved 79.5\% on AIME 2024, validating the effectiveness of this training methodology.

\paragraph{DAPO and the RLVR paradigm.}
\landmark{yu2025dapo} systematically improved upon GRPO,
boosting its performance on AIME from \textbf{30\% to 50\%} (7B model),
while providing a reproducible open-source practice for the
RLVR (Reinforcement Learning with Verifiable Rewards) paradigm.
DAPO's core technical contributions include four improvements:
\begin{enumerate}[leftmargin=2em]
  \item \textbf{Clip-higher}: In the clipping mechanism of PPO/GRPO,
    using a more relaxed upper bound clip ($1 + \varepsilon_{\text{high}}$) for positive advantage samples,
    encouraging the model to more boldly explore high-reward reasoning paths;
  \item \textbf{Dynamic sampling}: Dynamically adjusting the sampling strategy based on training progress,
    using higher temperature in early training to increase exploration, and gradually cooling down for exploitation later;
  \item \textbf{Token-level loss}: Decomposing GRPO's sequence-level loss
    into token-level loss, providing more fine-grained gradient signals;
  \item \textbf{KL constraint removal}: Completely removing the KL divergence penalty term,
    allowing the policy to diverge more freely from the reference model;
    on reasoning tasks, this ``unconstrained'' strategy yields significant gains.
\end{enumerate}
DAPO's success validates the core insight of the RLVR paradigm:
\textbf{for tasks with automated verifiers (mathematics, coding),
rule-based rewards are far more reliable and efficient than learned reward models}.
This conclusion has profoundly influenced the research direction of Thread~2 (\crossref{2}{sec:t2}).

\paragraph{The explosion of open-source reasoning models.}
2025 witnessed a Cambrian explosion of open-source reasoning models:
\begin{itemize}[leftmargin=2em]
  \item \textbf{Marco-o1 v2} \cite{yin2025marcoo1v2}:
    Explored the bottleneck problem of reasoning capability distillation
    and proposed methodologies for extending the distillation ceiling;
  \item \textbf{OpenR} \cite{wang2024openr}:
    Provided an open-source advanced reasoning framework
    integrating PRM, MCTS, and multiple RL algorithms;
  \item A large number of small reasoning models trained on R1 distillation data emerged,
    making reasoning capabilities no longer the exclusive province of large models.
\end{itemize}

\paragraph{``Thinking Models'' as the default paradigm.}
By early 2026, a clear industry consensus had formed:
\textbf{reasoning is no longer a separate model category, but a switchable model capability}.
The latest models from major providers including OpenAI, DeepSeek, Qwen, and Anthropic
all incorporate reasoning as a default or optional mode built into general-purpose models.
The significance of this shift is that
the goal of post-training has evolved from ``training a specialized model that can reason''
to ``training a reasoning mode that can be activated on demand within a general-purpose model.''

\subsubsection{Benchmarks and Methods for Mathematical Reasoning}
\label{sec:t4-math}

Mathematical reasoning is the most rigorous testing ground for verifying reasoning capabilities.
This section provides an overview of the key benchmarks and methods in this area.

\paragraph{Benchmark datasets.}
\begin{itemize}[leftmargin=2em]
  \item \textbf{MATH} \landmark{hendrycks2021math}:
    12,500 competition-level mathematics problems spanning 7 topics and 5 difficulty levels,
    each accompanied by a step-by-step solution.
    When released in 2021, the strongest model achieved only 6.9\% accuracy;
    by 2025, DeepSeek-R1 had reached 97.3\% (on the MATH-500 subset).

  \item \textbf{GSM8K} \cite{cobbe2021gsm8k}:
    8.5K grade school-level word problems requiring 2--8 reasoning steps.
    As an introductory benchmark, progress on GSM8K clearly reflects the evolution of reasoning methodologies.
\end{itemize}

\paragraph{SFT methods for mathematical reasoning.}
\begin{itemize}[leftmargin=2em]
  \item \textbf{WizardMath} \cite{luo2023wizardmath}:
    Used the Evol-Instruct strategy to generate diverse instruction data in the mathematics domain,
    performing SFT + RLHF on Llama-2 70B to achieve 81.6\% on GSM8K.

  \item \textbf{Llemma} \cite{azerbayev2023llemma}:
    Performed continued pre-training on Code Llama using a large-scale mathematical corpus (Proof-Pile-2, 55B tokens)
    to build a mathematics-oriented open language model.
    Llemma surpasses general-purpose models on multiple mathematical benchmarks even in the few-shot setting.

  \item \textbf{MAmmoTH} \cite{yue2023mammoth}:
    Constructed the MathInstruct dataset, mixing CoT and program-of-thought (PoT) reasoning approaches,
    using PoT for problems requiring precise computation and CoT for problems requiring conceptual reasoning.

  \item \textbf{OpenMathInstruct-1} \cite{toshniwal2024openmath}:
    Used Mixtral to generate mathematical reasoning data at scale, with rigorous answer-based filtering to ensure quality,
    constructing an open-source mathematical instruction dataset of 1.8M samples.

  \item \textbf{DART-Math} \cite{tong2024dartmath}:
    Proposed \textbf{Difficulty-Aware Rejection Tuning},
    which allocates a higher sampling budget to difficult problems during rejection sampling,
    addressing the issue of insufficient samples for hard problems in standard rejection sampling.
\end{itemize}

\subsubsection{Open-Source Reasoning Models}

The success of o1 and R1 sparked a wave of enthusiasm in the open-source community for replicating and innovating on reasoning models.

\begin{itemize}[leftmargin=2em]
  \item \textbf{Qwen2.5} \cite{qwen2024qwen25}:
    Alibaba's Qwen series substantially strengthened reasoning capabilities in version 2.5,
    achieving leading open-source model performance on mathematics and coding tasks through a carefully designed SFT + RL pipeline.

  \item \textbf{Sky-T1} \cite{novasky2025skyt1}:
    NovaSky-AI demonstrated that o1-level reasoning capabilities can be replicated
    with only \textbf{\$450} in training costs
    (approaching o1-preview on MATH-500 and AIME).
    This work dramatically lowered the barrier to entry for reasoning model research.

  \item \textbf{rStar} \cite{qi2024rstar}:
    Proposed \textbf{Mutual Reasoning}, having two smaller LLMs verify each other's reasoning processes through MCTS,
    improving small model reasoning capabilities without relying on large models.
    On GSM8K, this raised the accuracy of Phi3-mini 3.8B from 86.4\% to \textbf{90.0\%}.
\end{itemize}


% ============================================================
% 4.5 Test-Time Compute Scaling & Open Problems
% ============================================================
\subsection{Test-Time Compute Scaling and Open Problems}
\label{sec:t4-open}

\subsubsection{Test-Time Compute Scaling: A New Type of Scaling Law}
\label{sec:t4-ttc}

Traditional LLM scaling laws \cite{kaplan2020scaling,hoffmann2022chinchilla}
focus on \textbf{training-time} resource allocation:
given a fixed training compute budget, how to optimally allocate between model size and training data volume.
In 2024, a complementary scaling paradigm emerged:
\textbf{test-time compute scaling}, which continuously improves performance by investing more computation during the inference stage.

\paragraph{Snell et al.'s compute-optimal framework.}
\landmark{snell2024testtimescaling} was the first to systematically study
the optimal allocation strategy for test-time compute.

Snell et al.\ considered two mechanisms for leveraging additional computation at inference time:
\begin{enumerate}[leftmargin=2em]
  \item \textbf{Search against a verifier}:
    Generating multiple candidate responses and using a PRM for search and ranking.
    This includes strategies such as best-of-$N$ sampling and beam search.

  \item \textbf{Revising proposals}:
    Having the model iteratively revise its response based on previous attempts.
    By providing previous (incorrect) responses as context,
    the model generates improved versions.
\end{enumerate}

\paragraph{Core finding: Difficulty-dependent optimal strategies.}
The most critical finding is that \textbf{the optimal allocation strategy for test-time compute depends on the difficulty of the problem}.
\begin{itemize}[leftmargin=2em]
  \item \textbf{Easy problems}:
    Revision-based strategies (iterative modification) are more effective,
    as the model only needs to make minor corrections on an already decent answer;
  \item \textbf{Hard problems}:
    Search-based strategies (parallel sampling + PRM verification) are more effective,
    as hard problems require exploring a broader space of solution paths.
\end{itemize}

\noindent
Based on this finding, Snell et al.\ proposed a \textbf{compute-optimal} strategy:
dynamically selecting the best test-time compute allocation method based on problem difficulty.
Under the same test-time compute budget,
the compute-optimal strategy saves \textbf{4$\times$} computational resources compared to best-of-$N$.

\paragraph{Smaller model + more thinking time $\approx$ larger model.}
The most practically significant finding is that
on the MATH benchmark,
\textbf{a smaller model with sufficient test-time compute
can match or even exceed the performance of a model 14$\times$ larger in parameters}.
This means that in certain deployment scenarios,
a small and flexible model with a dynamic reasoning budget can replace a large and cumbersome model,
achieving a better cost-performance tradeoff.

\paragraph{Inference Scaling Laws.}
Wu et al.\ \cite{wu2024inferencescaling} studied inference scaling from a more theoretical perspective.
They found that the gains from test-time compute follow a \textbf{log-linear} pattern:
each doubling of test-time compute yields roughly constant but diminishing performance improvement,
similar to the power law observed in training-time scaling.

\paragraph{s1: Simple budget forcing.}
Muennighoff et al.\ \cite{muennighoff2025s1} proposed
an extremely concise test-time scaling method called \textbf{s1}:
during reasoning, special tokens are inserted to ``force'' the model to continue thinking (budget forcing),
or an end-of-thinking token is inserted early to shorten reasoning.
Using only 1,000 curated examples + budget forcing,
s1 achieved performance on AIME 2024 \textbf{comparable to o1-preview},
demonstrating the simplicity and effectiveness of test-time compute scaling methods.

\subsubsection{Unified Framework: Train-Time vs Test-Time Compute}

Viewing train-time and test-time compute scaling in a unified manner,
Zeng et al.\ \cite{zeng2024roadmapo1} provided a systematic roadmap:
\begin{itemize}[leftmargin=2em]
  \item \textbf{Train-time scaling} (larger models, more data)
    $\rightarrow$ stronger base capability;
  \item \textbf{Test-time scaling} (more sampling, deeper search, longer reasoning chains)
    $\rightarrow$ more thorough utilization of existing capabilities;
  \item \textbf{Post-training} (RL, PRM, GRPO)
    $\rightarrow$ the bridge connecting the two---the reasoning strategies acquired through post-training
    determine the efficiency with which test-time compute is utilized.
\end{itemize}

\subsubsection{Open Problems}

\begin{openproblembox}
\textbf{Open Problem 1: Automation and Generalization of Process Supervision}\\
PRM has demonstrated its value in mathematical reasoning,
but how to extend process supervision to \textbf{domains where formal verification is not possible}
(such as scientific reasoning, legal reasoning, and logical consistency in creative writing)
remains a fundamental open problem.
When there is no ground truth answer for a completer model to verify,
how can reliable step-level rewards be automatically generated?
\end{openproblembox}

\begin{openproblembox}
\textbf{Open Problem 2: The Scaling Ceiling of Reasoning}\\
DeepSeek-V3.2 \cite{deepseek2025v32} achieved gold medal level in IMO-level mathematics competitions,
and o3 reached 88.9\% on AIME 2025 \cite{openai2025o3},
indicating that traditional competition benchmarks are rapidly saturating.
But does this mean reasoning capabilities have been ``solved''?
On open-ended mathematical problems requiring creative insight (rather than standard solution techniques),
multi-step scientific derivations, and cross-domain reasoning,
where is the ceiling of current models?
More critically, there is a debate in the RLVR paradigm about whether it ``makes models faster not smarter''---
methods like DAPO significantly improve pass rates on known task types,
but do they truly expand the frontier of model reasoning capabilities?
\end{openproblembox}

\begin{openproblembox}
\textbf{Open Problem 3: The Efficiency Frontier of Test-Time Compute}\\
Current test-time compute scaling methods
(whether parallel sampling or sequential revision)
exhibit logarithmically diminishing efficiency gains.
Is there a fundamentally more efficient way to utilize test-time compute?
For example, can models learn to \textbf{adaptively} allocate their reasoning budget---
answering easy problems quickly while investing substantial thought in hard problems---
without requiring an external difficulty estimator?
\end{openproblembox}

\begin{openproblembox}
\textbf{Open Problem 4: Theoretical Understanding of RL for Reasoning}\\
The emergent reasoning phenomena observed in DeepSeek-R1-Zero
(self-reflection, aha moments, automatic lengthening of reasoning chains)
lack a theoretical explanation.
Why can simple binary rewards + GRPO give rise to such complex reasoning behaviors?
Does this emergence require specific pre-training conditions (e.g., the training data containing sufficient reasoning corpora)?
Can it be reliably reproduced on smaller models?
\end{openproblembox}

\begin{openproblembox}
\textbf{Open Problem 5: Reasoning Reliability and Safety}\\
Long-chain reasoning introduces new safety challenges:
hidden CoT may contain inappropriate content (e.g., ``reasoning'' about harmful topics),
the opacity of reasoning chains makes auditing difficult
(\crossref{3}{sec:t3}),
and unbounded expansion of test-time compute could enable denial-of-service (DoS) attacks.
How can the safety and auditability of the reasoning process be ensured while maintaining reasoning capabilities?
\end{openproblembox}

\begin{openproblembox}
\textbf{Open Problem 6: Optimal Strategies for Thinking Mode Fusion}\\
Models such as DeepSeek V3.1 and Qwen3 \cite{yang2025qwen3} have implemented
the fusion of reasoning and non-reasoning modes within a single model.
However, current fusion strategies remain coarse---
the decision boundary for when the model should ``think deeply'' versus ``answer quickly'' is still unclear.
The more fundamental question is:
how to optimize this mode-switching strategy during the post-training stage
so that the model achieves Pareto optimality between reasoning depth and response latency?
Qwen3's four-stage pipeline provides one explicit training approach,
but the optimal mixing ratios and training strategies for the third-stage thinking mode fusion
still lack systematic investigation.
\end{openproblembox}

\subsection{Chapter Summary}

\begin{figure}[!ht]
\centering
\begin{tikzpicture}[
  node distance=0.6cm and 1.2cm,
  every node/.style={font=\small},
  milestone/.style={rectangle, rounded corners, draw=thread4, fill=thread4!10,
    text width=4.2cm, minimum height=0.8cm, align=center, font=\small\bfseries},
  branch/.style={rectangle, rounded corners, draw=gray!60, fill=gray!5,
    text width=3.2cm, minimum height=0.7cm, align=center, font=\small},
  arrow/.style={-{Stealth[length=2.5mm]}, thick, thread4!70},
  brancharrow/.style={-{Stealth[length=2mm]}, gray!60, thick},
]
% Main timeline
\node[milestone] (cot) {CoT (Wei 2022)\\Chain-of-Thought Prompting};
\node[milestone, below=1.2cm of cot] (star) {STaR (Zelikman 2022)\\Self-Taught Reasoner};
\node[milestone, below=1.2cm of star] (prm) {PRM (Lightman 2023)\\Process Reward Model};

% Branches
\node[branch, below left=1.2cm and 0.5cm of prm] (brA) {Branch A: Inference-Time\\Scaling: Self-Consistency,\\Best-of-N};
\node[branch, below right=1.2cm and 0.5cm of prm] (brB) {Branch B: RL Training\\GRPO, REINFORCE++};

% Convergence
\node[milestone, below=2.8cm of prm] (o1) {o1 (OpenAI 2024)\\Thinking Model Paradigm};
\node[milestone, below=1.2cm of o1] (r1) {DeepSeek-R1 (2025)\\RL from Scratch};
\node[branch, below=1.2cm of r1] (conv) {Convergence: Thinking\\Models as Industry Default};

% Arrows
\draw[arrow] (cot) -- (star);
\draw[arrow] (star) -- (prm);
\draw[brancharrow] (prm) -- (brA);
\draw[brancharrow] (prm) -- (brB);
\draw[arrow] (prm) -- (o1);
\draw[arrow] (o1) -- (r1);
\draw[brancharrow] (brA) -- (o1);
\draw[brancharrow] (brB) -- (o1);
\draw[arrow] (r1) -- (conv);
\end{tikzpicture}
\caption{Evolutionary path of Thread 4: from chain-of-thought prompting to thinking models as the default paradigm.
Purple nodes represent landmark works; gray nodes represent important branches.}
\label{fig:reasoning-evolution}
\end{figure}

\begin{table}[!ht]
\centering
\caption{Comparison of core methods in Thread 4. Training Signal distinguishes prompt-time, inference-time, and training-time methods.}
\label{tab:reasoning-comparison}
\small
\begin{tabularx}{\textwidth}{l >{\raggedright\arraybackslash}p{2.8cm} >{\raggedright\arraybackslash}X >{\raggedright\arraybackslash}X}
\toprule
\textbf{Method} & \textbf{Training Signal} & \textbf{Verification} & \textbf{Key Result} \\
\midrule
CoT Prompting & None (prompt-time) & N/A & $+$15\% on GSM8K (540B PaLM) \\
Self-Consistency & None (inference-time) & Majority vote over $K$ samples & $+$17.9\% on GSM8K over greedy CoT \\
STaR & Self-generated rationales & Answer correctness filter & Bootstraps reasoning from few examples \\
PRM & 800K step-level labels & Per-step human correctness & 78.2\% MATH (vs 72.4\% ORM) \\
GRPO & Group-normalized advantage & Rule-based (accuracy) & 88.2\% GSM8K, 51.7\% MATH \\
o1 & Large-scale RL & Hidden CoT + test-time scaling & 83.3\% AIME (vs 13.4\% GPT-4o) \\
R1-Zero & GRPO, binary reward only & Accuracy only; no SFT & 71.0\% AIME; emergent self-reflection \\
R1 & {\scriptsize SFT$\to$RL$\to$RS$\to$RL} & Rule + preference RM & 79.8\% AIME; distillable \\
\bottomrule
\end{tabularx}
\end{table}

The evolutionary trajectory of Thread~4 clearly demonstrates a core paradigm shift in the post-training landscape:
from \textbf{inference-time prompting} (CoT, Self-Consistency)
to \textbf{training-time internalization} (STaR, GRPO, R1),
to \textbf{deep integration of both} (test-time compute scaling with RL-trained models),
and finally to \textbf{reasoning as the default mode} (thinking models becoming the industry standard, 2025--2026).

\evolution{CoT Prompting (2022): reasoning as a prompt technique}%
{Thinking Models (2025--2026): reasoning as the model's default mode}%
{RLVR + Thinking Mode Fusion}

\noindent
Several core insights:
\begin{enumerate}[leftmargin=2em]
  \item \textbf{Data structures determine everything}:
    The transition from ORM to PRM is essentially an upgrade of the ``data structure'' for reasoning verification---
    from sequence-level labels to step-level labels.
    As Linus observed, good programmers focus on data structures, not code.

  \item \textbf{Simple methods + sufficient compute $>$ complex methods}:
    GRPO removes the value model from PPO,
    R1-Zero uses the simplest binary reward,
    and s1 uses only 1,000 examples + budget forcing.
    In the reasoning domain, \textbf{simple designs + sufficient train/test-time compute}
    have repeatedly outperformed complex engineering solutions.

  \item \textbf{A paradigmatic example of cross-thread idea flow}:
    Thread~4 is perhaps the thread with the most intensive interactions with other threads:
    reward modeling flows from Thread~2, safety considerations from Thread~3,
    downstream applications flow to Thread~6, and efficiency requirements come from Thread~7.
    The improvement of reasoning capabilities is the result of co-evolution across the entire post-training ecosystem.
\end{enumerate}

\paragraph{Preview of the next chapter.}
The reasoning capabilities of thinking models are currently demonstrated primarily in the text domain.
When reasoning extends to the visual domain, models need to first ``see'' before they can ``reason.''
Thread~5 will trace the evolution of multimodal post-training,
from the proposal of visual instruction tuning to hallucination mitigation,
and then to the integration of multimodal reasoning---showing how the post-training methodology for LLMs transfers to and evolves within VLMs.
